\section{Dimensions}

Regularity behaves differently in each dimension, and the base lives where it behaves most richly. Climb the ladder one rung at a time. In two dimensions the regular figures are the polygons, and there are infinitely many, one for every number of sides. In three dimensions there are exactly five, the Platonic solids. In four there are six. In five and in every dimension above, there are only three, the simplex, the cube, and the cube's dual. The variety peaks at four and then collapses.

\begin{table}[ht]
\centering
\begin{tabular}{@{}rlp{0.4\textwidth}@{}}
\toprule
dimension & count & the regular figures \\
\midrule
$2$ & $\infty$ & the polygons, one per side count \\
$3$ & $5$ & tetrahedron, cube, octahedron, \newline dodecahedron, icosahedron \\
$4$ & $6$ & 5-cell, 8-cell, 16-cell, 24-cell, \newline 120-cell, 600-cell \\
$\geq 5$ & $3$ & simplex, hypercube, cross-polytope \\
\bottomrule
\end{tabular}
\caption{Regular polytopes by dimension. Four is the richest, and the only one with a figure that has no analog above or below.}
\label{tab:regular}
\end{table}

Four dimensions holds a figure that exists nowhere else, the 24-cell. Five of the six four-dimensional regulars are higher versions of three-dimensional solids, but the 24-cell is a version of nothing, with no counterpart in three dimensions or in any other. The base is built from it. The next section shows why that uniqueness matters. For now it is enough that the one dimension with a regular figure that stands alone is the one the bulk lives in.

The bulk is four-dimensional, and its flat boundary, the cusp we live in, is three-dimensional. Neither number is assumed. Both fall out of the selection argument a few sections on, where the demand for a clean three-dimensional space of ordinary physics, together with a dock that carries spin, forces the bulk up to four and pins the cusp at three. The ladder itself runs on past four, with tessellations and curvatures in every dimension up to infinity, all governed by the same three-sign rule. But the regular structures thin to almost nothing above four, and nothing up there offers what the 24-cell offers, so the story stays in three and four.

Below the committed bulk runs a whole ladder of hyperbolic substrates, one per dimension, and they are worth seeing together. In two dimensions the heptagonal tiling, in three the dodecahedral honeycomb $\{5,3,4\}$, in four the committed $\{3,4,3,4\}$, and on up to five, each a different object with its own cells, its own symmetry, and its own spin structure. They are not copies of one another, and no map turns one into the next for free. Going down a dimension is exact but lossy, the cells of a honeycomb are the honeycomb one dimension lower, read off by dropping the last entry of its symbol, and its boundary or its flat horosphere is lower still, but each forgets the bulk it came from. Going up is possible and never free, by holography, which rebuilds a bulk from its boundary, or by giving each cell an internal thread that becomes a new direction, and the added dimension's information must be supplied, never conjured from the bare lower graph. The ladder is a family joined by maps, like the triangle, the tetrahedron, and the 5-cell, each carrying the last as a face and none equal to another.

What rides the whole ladder is the spinor. Every rung carries one, in the double cover of its own rotation symmetry, the order-$120$ icosians on the three-dimensional rung, the order-$24$ binary tetrahedral group on the four-dimensional one. And these covers are not unrelated. By the McKay correspondence, a deep dictionary between the finite rotation groups and the exceptional Lie algebras, the binary tetrahedral group answers to $E_6$, the binary octahedral to $E_7$, and the binary icosahedral to $E_8$, so the spin structures of the substrates across dimensions are three faces of the one exceptional tower $E_6 \subset E_7 \subset E_8$. The dock's own group sits at the foot of that tower, one more sign that the substrate is assembled from the most exceptional objects mathematics has.
